\setmainhead{\emph{\textsc{Ilíada}}, Canto VIII}
\bnum
\section*{\centering CANTO VIII\footT{\ehdos, pp. 102-104.}}
\addcontentsline{toc}{subsection}{\textsc{Canto VIII}}
%
\separator
\postart
\setotherline{130}
Y allá entre mil estragos el desastre ocurriera,\\
y los de Ilión quedaran cual corderos cercados,\\
si al punto el padre de hombres y dioses no lo viera.\\
Tronando, pues, horrísono, lanzó a tierra el fulgente\\
rayo ante los corceles de Diomedes que, entonces,\\
al surgir la terrible llama de azufre ardiente,\\
recularon de espanto bajo el carro de bronce;\\
y escapándose a Néstor el rendaje luciente,\\
\dlog{con miedo en su alma díjole a Diomedes:}{---¡Tidida,}\\
Revolvamos los fuertes potros para escaparnos!\\
¿No ves que el triunfo niégate Dios, y que hoy Zeus Kronida\\
otorga a ése la gloria que otra vez querrá darnos?\\
Pues por bravo que sea, no habrá varón que acierte\\
a estorbar el designio de Dios, mucho más fuerte.\par
Diomedes el cumplido guerrero, así le dijo:\\
---Sí, anciano, lo que hablaste, justo y cuerdo es, de fijo;\\
mas, pecho y alma embárgame cruel dolor, porque un día,\\
dirigiendo Héctor a los troyanos la palabra,\\
dirá que de él Tideides hacia la flota huía.\\
¡Y antes que así se alabe, que la tierra se me abra!\par
Mas, respondióle Néstor Gerenio el caballero:\\
---¡Ay de mí, qué has dicho hijo de Tideo el guerrero!\\
Aunque Héctor vil y flojo te llame, tales cosas\\
ni troyanos ni dárdanos creerán, ni las esposas\\
troyanas cuyos bravos y equipados maridos,\\
tú, en la flor de los años, entre el polvo has tendido.\par
Así habló, y revolviendo los corceles gallardos,\\
entre el tumulto huyeron; mas, con ruido inquietante,\\
hacían los troyanos y Héctor llover sus dardos;\\
y alzó el grito el grande Héctor del casco tremolante:\par
---¡Tideides, si los dánaos de la caballería,\\
grande honra con el solio, mesa y brindis te hacían,\\
van a despreciarte ahora que como mujer corres!\\
huye, pues, vil muchacha! Ni escalar nuestras torres\\
te permitiré, ni hacia los navíos llevarte\\
las mujeres, pues previo castigo pienso darte!\par
Dijo, y dudó Tideides entre seguir la huída,\\
o tornar los caballos para pelear de frente.\\
Tres veces vacilaron su corazón y mente,\\
y otras tantas el próvido Zeus, desde el monte Ida,\\
tronó por los troyanos, en signo fehaciente\\
de la victoria al bélico azar indefinida.\par
\poend
\separator
%
\postart
\setotherline{191}
¡Seguidme al punto, y rápidos, veamos de arrancar\\
su escudo a Néstor, célebre en la propia mansión\\
celeste, porque es todo de oro, hasta la armazón;\\
y al jinete Diomedes de los hombros sacar\\
el rico arnés que Hefesto le cinceló; pues creo\\
que si logramos esas dos cosas, los aqueos\\
esta misma noche hácense en su flota a la mar!\par
\poend
\separator
%
\postart
\setotherline{253}
No hubo quien de los muchos dánaos aventajara\\
al Tideida, en lanzarse con los raudos corceles\\
fuera del foso, a objeto de pelear cara a cara.\par
\poend
\separator
%
\postart
\setotherline{492}
Y aquéllos desmontaron para escuchar entonces\\
lo que les decía Héctor caro a Zeus, empuñando\\
su lanza de once codos cuya punta de bronce\\
que un aro de oro afianza, se erige relumbrando.\\
Y así, apoyado en ella, los fué aquél arengando.\par
\poend
\separator
%
\postart
\setotherline{529}
Hagamos, pues, la guardia nocturna, y bien temprano\\
salgamos pertrechados de todas armas, y ante\\
la flota, al ardiente Ares convoquemos ufanos.\\
Y sabré si Diomedes, el Tideida pujante,\\
desde la escuadra al muro me arroja, o si lo mato\\
con el bronce, y sangrientos despojos le arrebato.\\
Mañana ha de probarnos su valor con su aguante\\
cuando lo ataque a lanza; mas, creo que adelante\\
caérá herido junto con muchos camaradas\\
al salir el sol.\par
\poend
\separator
\enum